%! Unit 1: Vectors and Matrices — Summative Assessment — Unit Test (blank)
\documentclass[10pt]{article}
\usepackage{linalg-article}
\usepackage{linalg-boxes}
\newcommand{\vv}{\mathbf{v}}
\newcommand{\ww}{\mathbf{w}}
\newcommand{\uu}{\mathbf{u}}
\newcommand{\bb}{\mathbf{b}}
\newcommand{\xx}{\mathbf{x}}

% Part-divider strip
\newcommand{\parthead}[1]{%
  \vspace{6pt}\noindent
  \begin{headlinebox}{burgundy}{\color{white}\bfseries #1}\end{headlinebox}%
  \vspace{4pt}}

\begin{document}

\pageheader{Unit 1: Vectors and Matrices}{Unit Test}
\namedateperiod

\vspace{0.06in}
\noindent\textbf{Instructions:} Show all work. No notes unless your teacher says so.

% ======================================================================
\parthead{Part A --- Vocabulary (8 pts)}
Write the letter of the best description next to each term.

\vspace{4pt}
\noindent
\begin{minipage}[t]{0.44\textwidth}
\begin{enumerate}[leftmargin=2.2em, itemsep=7pt, label=\arabic*.]
  \item Linear combination \blank{1.2cm}
  \item Rank \blank{1.2cm}
  \item Dot product \blank{1.2cm}
  \item Unit vector \blank{1.2cm}
  \item Column space \blank{1.2cm}
  \item Scalar \blank{1.2cm}
  \item Span \blank{1.2cm}
  \item Magnitude \blank{1.2cm}
\end{enumerate}
\end{minipage}%
\hfill
\begin{minipage}[t]{0.53\textwidth}
\small
\begin{description}[leftmargin=1.4em, itemsep=3pt]
  \item[(A)] The length of a vector, $\sqrt{v_1^2+v_2^2}$.
  \item[(B)] The set of \emph{all} linear combinations of a set of vectors.
  \item[(C)] A single number that scales (stretches or flips) a vector.
  \item[(D)] All vectors you can reach as $A\xx$ --- the span of a matrix's columns.
  \item[(E)] A vector whose length is exactly $1$.
  \item[(F)] $\vv\cdot\ww=v_1w_1+v_2w_2$; it measures length and angle.
  \item[(G)] The number of independent columns of a matrix.
  \item[(H)] A sum of scaled vectors, $c\vv+d\ww$.
\end{description}
\end{minipage}

% ======================================================================
\parthead{Part B --- Multiple Choice (12 pts)}
Circle the best answer.

\begin{enumerate}[leftmargin=1.9em, itemsep=9pt]
  \item Two nonzero vectors are \textbf{perpendicular} exactly when
    \hfill (A) $\vv\cdot\ww=1$ \quad (B) they point the same way \quad
    (C) $\vv\cdot\ww=0$ \quad (D) they have equal length
  \item The vectors $\begin{bsmallmatrix}2\\1\end{bsmallmatrix}$ and
    $\begin{bsmallmatrix}4\\2\end{bsmallmatrix}$ span
    \hfill (A) the whole plane \quad (B) a line \quad (C) a single point \quad (D) all of space
  \item In the product $A\xx$, the entries of $\xx$ are
    \hfill (A) the lengths of the rows \quad (B) the weights on the \emph{columns} of $A$ \\
    \phantom{x}\hfill (C) always $0$ or $1$ \quad (D) the angle between columns
  \item The matrix $\begin{bsmallmatrix}3&6\\1&2\end{bsmallmatrix}$ has rank
    \hfill (A) $0$ \quad (B) $1$ \quad (C) $2$ \quad (D) $6$
  \item If $\vv$ has length $10$, then the unit vector $\dfrac{\vv}{\|\vv\|}$ has length
    \hfill (A) $10$ \quad (B) $100$ \quad (C) $1$ \quad (D) $\tfrac{1}{10}$
  \item Every vector $\bb$ in the plane is reachable as $A\xx$ exactly when $A$'s two columns are
    \hfill (A) parallel \quad (B) perpendicular \quad (C) unit vectors \quad (D) independent
\end{enumerate}

% ======================================================================
\parthead{Part C --- Short Answer \& Computation (30 pts)}
Show your work.

\begin{enumerate}[leftmargin=1.9em, itemsep=13pt]
  \item Let $\vv=\begin{bsmallmatrix}5\\12\end{bsmallmatrix}$ and
        $\ww=\begin{bsmallmatrix}-2\\3\end{bsmallmatrix}$. Compute
        \textbf{(a)} $\vv+\ww$, \quad \textbf{(b)} $2\vv-\ww$, \quad \textbf{(c)} $\|\vv\|$.
        \vspace{1.6cm}
  \item Write the unit vector that points the same direction as
        $\vv=\begin{bsmallmatrix}5\\12\end{bsmallmatrix}$.
        \vspace{1.2cm}
  \item Let $\uu=\begin{bsmallmatrix}3\\6\end{bsmallmatrix}$ and
        $\vv=\begin{bsmallmatrix}2\\-1\end{bsmallmatrix}$. Compute $\uu\cdot\vv$, then state
        whether $\uu$ and $\vv$ are perpendicular and how you know.
        \vspace{1.6cm}
  \item Find weights $c$ and $d$ so that
        $c\begin{bsmallmatrix}3\\1\end{bsmallmatrix}+d\begin{bsmallmatrix}1\\2\end{bsmallmatrix}
         =\begin{bsmallmatrix}9\\8\end{bsmallmatrix}$.
        \vspace{1.6cm}
  \item Let $A=\begin{bsmallmatrix}2&1\\1&3\end{bsmallmatrix}$ and
        $\xx=\begin{bsmallmatrix}3\\1\end{bsmallmatrix}$. Compute $A\xx$ \emph{as a combination
        of the columns of $A$}.
        \vspace{1.6cm}
  \item Factor $A=\begin{bsmallmatrix}1&2&4\\1&1&3\end{bsmallmatrix}$ as $A=CR$. Which columns are
        independent? Write $C$, build the recipe $R$, and give the rank.
        \vspace{1.8cm}
  \item Let $A=\begin{bsmallmatrix}1&3\\2&6\end{bsmallmatrix}$. For each target, say whether it is
        reachable as $A\xx$ (is it in the column space?) and justify in one line.
        \textbf{(a)} $\bb=\begin{bsmallmatrix}2\\4\end{bsmallmatrix}$ \quad
        \textbf{(b)} $\bb=\begin{bsmallmatrix}3\\3\end{bsmallmatrix}$
        \vspace{1.6cm}
\end{enumerate}

% ======================================================================
\parthead{Part D --- Extended Response (10 pts)}
Explain your reasoning in full sentences.

\begin{enumerate}[leftmargin=1.9em, itemsep=16pt]
  \item A matrix $A=\begin{bsmallmatrix}2&1\\1&1\end{bsmallmatrix}$ has columns
        $\begin{bsmallmatrix}2\\1\end{bsmallmatrix}$ and $\begin{bsmallmatrix}1\\1\end{bsmallmatrix}$.
        \textbf{(a)} Is the column space a line or the whole plane? How do you know?
        \textbf{(b)} Is $\bb=\begin{bsmallmatrix}5\\3\end{bsmallmatrix}$ reachable? If so, find the
        weights. \textbf{(c)} If the second column were changed to
        $\begin{bsmallmatrix}4\\2\end{bsmallmatrix}$, what would happen to the column space, and why?
        \vspace{2.4cm}
  \item Two shoppers' preference profiles are $\uu=\begin{bsmallmatrix}3\\2\end{bsmallmatrix}$ and
        $\vv=\begin{bsmallmatrix}-2\\3\end{bsmallmatrix}$. Compute $\uu\cdot\vv$ and explain what the
        result says about how similar the two shoppers' preferences are, using what you know about
        the sign of a dot product.
        \vspace{2.0cm}
\end{enumerate}

\end{document}
